For each $t\geq1$ and $\epsilon>0$, must every
$A\subseteq\{1,\ldots,t\}^N$ of size at least $\epsilon t^N$ contain
a combinatorial line when $N$ is sufficiently large?
A line has a nonempty set $S\subseteq\{1,\ldots,N\}$ of variable
coordinates: its $i$th word has symbol $i$ in every coordinate of $S$
and a fixed common symbol at each remaining coordinate.

The general assertion is the established density Hales--Jewett theorem
of Furstenberg and Katznelson, with a later combinatorial proof by the
Polymath project. This attempt gives a complete elementary proof for
$t\leq2$. It does not reconstruct the theorem for all alphabets.
